% letter-academic-two-page.tex — layout stress: an academic letter with enough
% prose to span two pages.
%
% This is the only fixture that can exercise the academic letter's running
% header, which appears from page two onwards. letter-academic-long.tex stresses
% long recipient fields but fits on one page, and letter-two-page.tex is an
% industry letter, whose empty page style has no furniture at all — so without
% this fixture the running-header assertion in run.sh would never execute.
%
% Verifies that the page break falls inside the body, that page two carries the
% `<name> -- Cover Letter` running header, and that every page carries a
% centered `Page N of M` folio matching careerdossier-cv.
\documentclass[family=academic]{careerdossier-letter}
\CDossierSetup{
  name     = {Ada Lovelace},
  headline = {Researcher in Analytical Computing},
  email    = {ada.lovelace@example.com},
  location = {Toronto, Ontario, Canada},
  scholar  = {scholar.google.com/citations?user=ada-lovelace}
}
\CDossierLetterSetup{
  date                   = {17 July 2026},
  recipient-name         = {Professor Grace Hopper},
  recipient-title        = {Chair of Computational Science},
  recipient-organization = {Example University},
  recipient-address      = {100 Academic Avenue\\Toronto, Ontario M5S 1A1},
  subject                = {Application for the Assistant Professor Position},
  salutation             = {Dear Professor Hopper,}
}
\begin{document}
\MakeCDossierLetterhead

I am writing to apply for the Assistant Professor position in Computational
Science. My research connects reliable numerical methods, reproducible software,
and accessible teaching materials for interdisciplinary scientific work.

My recent work has focused on creating computational methods that can be
understood, tested, and reused by collaborators across mathematics, engineering,
and data science. I develop clear research software, publish durable technical
documentation, and involve students in the full process from modelling through
validation and communication.

Reproducibility is the organizing principle of my research programme. Every
result I publish is accompanied by the software, data preparation steps, and
environment description needed to regenerate it from source. This practice began
as a personal discipline and has become a collaborative one: my students learn to
treat an unreproducible result as an incomplete result, and reviewers of my work
can inspect the computational path from raw input to published figure.

My teaching follows the same commitment. I have designed and delivered courses in
numerical analysis, scientific computing, and statistical methods, each built
around the idea that students should be able to reconstruct every result they are
shown. Course materials are versioned, openly licensed, and written so that a
student returning to them years later can still run the examples. I have
supervised undergraduate and graduate projects spanning method development,
software engineering practice, and the communication of quantitative results to
non-specialist audiences.

Beyond research and teaching, I have contributed to the maintenance of open
research software used across several institutions, served on programme
committees for conferences concerned with reproducible computational science, and
organized workshops introducing early-career researchers to version control,
automated testing, and durable documentation practices.

I would welcome the opportunity to contribute to your department's research,
teaching, and student mentorship. I am especially interested in building courses
that connect theoretical foundations with transparent computational practice, and
in collaborating with colleagues whose scientific questions would benefit from
carefully engineered numerical methods.

\MakeCDossierClosing
\end{document}
